\honest{A Rational Argument}{Eliezer Yudkowsky}{2007}

You are a campaign manager, just hired by Mortimer Q. Snodgrass, the Green candidate for Mayor of Hadleyburg. \nb{Hadleyburg is the town of Mark Twain's story ``The Man That Corrupted Hadleyburg'' (1899), which is proud of an honesty that fails its first real test.} You want ``an impeccable rational argument'' that he is the best candidate. ``Sorry. It can't be done.'' What if every fact you cite is true and relevant? ``Sorry. It still can't be done. You defeated yourself the instant you specified your argument's conclusion in advance.'' \nb{That is true of arguments that pick among evidence, where what is left out changes what the rest shows. It is not true of deduction: a proof is often written toward a conclusion known in advance, and it is valid or not whatever the order of writing. The post's own last paragraph defines a ``rational'' argument as one whose conclusion was written after its premises, and in that sense the verdict holds by definition.}

The local newspaper sent the candidates a 16-question questionnaire, and then its offices were destroyed by a meteorite. Your candidate compares well on 15 questions. The exception is question 11: ``Are you now, or have you ever been, a supervillain?'' You are tempted to publish the questionnaire without question 11. That crosses the line from rationality to rationalization: voters must now condition on the fact of your presentation and infer hidden evidence. In fact you crossed the line when you checked whether the questionnaire helped before deciding to publish it. A voter still choosing a candidate would not censor useful information.

A ``logical'' argument is one that follows from its premises. ``All rectangles are quadrilaterals. All squares are quadrilaterals. Therefore, all squares are rectangles'' is illogical, though everything in it is true. It is worth refusing to excuse such a deduction even when its conclusion is true. The difference may matter when new evidence comes in, and sloppiness is habit-forming. Above all, it gives the wrong explanation. Squares may be rectangles, but not because both are quadrilaterals. Such a syllogism is ``hypocritical'': its stated reasons are not its real reasons. \nb{Its premises do not support its conclusion at all. The campaign's facts do support its conclusion; their fault is what was left out. The syllogism shows a different error from the questionnaire.}

The only honest method: before anyone hires you, gather the evidence, make your checklist, find the winner, offer to manage the campaign, and print your checklist as campaign literature. ``Whatever actually decides your bottom line is the only thing you can honestly write on the lines above.'' \nb{That is a rule about the honesty of the arguer, not about the facts cited. In a comment the same day Yudkowsky wrote: ``Valid evidence is valid, whatever the motives of the one who cites it.''}
